\section{Fundamentos y estado del arte}\label{sec:estado}

\subsection{Pronóstico fotovoltaico operacional, heterogeneidad y normalización física}

La brecha entre el rendimiento teórico de placa y la energía medida en centrales conectadas a red constituye un fenómeno bien documentado en la literatura empírica. Müller et~al.~\cite{mueller2016yield} contrastaron estimaciones de rendimiento frente a mediciones reales de $26$ plantas en Alemania y España, mostrando que los sesgos en bases de irradiación y las pérdidas operativas generan desviaciones sistemáticas; Purohit y Purohit~\cite{purohit2018performance} observaron discrepancias análogas en centrales de la India. Al extender el análisis desde instalaciones aisladas hacia flotas heterogéneas o modelos transferibles entre emplazamientos, Trigo-González et~al.~\cite{trigo2019hourly,trigo2021models} formularon regresiones lineales múltiples y modelos no lineales sobre el \textit{Performance Ratio} en Chile y España bajo pérdida cuadrática, mientras que Lehmann et~al.~\cite{lehmann2023benchmark} compararon ocho soluciones comerciales intradía. En el contexto colombiano, Pico et~al.~\cite{pico2026shortterm} evaluaron algoritmos de aprendizaje automático a resolución de diez minutos para el horizonte de un día sobre $7$ plantas; sin embargo, emplearon potencia simulada mediante SAM/NSRDB y partición puramente cronológica, sin generación real medida por el operador del mercado ni horizontes subestacionales.

Para comparar de manera coherente centrales cuya capacidad instalada difiere en órdenes de magnitud, Yang et~al.~\cite{yang2020} sistematizaron buenas prácticas de verificación determinística solar, recomendando normalizar las métricas por capacidad instalada o cielo despejado, reportar \textit{skill scores} frente a referencias causales ingenuas y preservar la no anticipación temporal al procesar salidas numéricas del tiempo (NWP). Esta recomendación motiva directamente la normalización causal de SM-CFHR por $24\,000\,C_{i,t}$: mientras la literatura suele normalizar únicamente las métricas finales de evaluación, en flotas con $39{,}6\times$ de diferencia de escala y ampliaciones escalonadas de potencia es igualmente necesario normalizar la variable objetivo y los predictores históricos antes de optimizar la función de pérdida.

\subsection{Pronósticos diarios, subestacionales y posprocesamiento meteorológico}

En horizontes cortos ($D+1$), los pronósticos NWP presentan sesgos locales sistemáticos debidos a topografía, nubosidad convectiva subgrilla y parametrización atmosférica. Glahn y Lowry~\cite{glahn1972mos} introdujeron el posprocesamiento estadístico (\textit{Model Output Statistics}, MOS) para corregir salidas numéricas mediante regresión sobre observaciones locales; posteriormente, Bacher et~al.~\cite{bacher2009online} acoplaron predicciones NWP con dinámica autorregresiva reciente en generación solar en línea, y Yang y van der Meer~\cite{yang2021postprocessing} sintetizaron los principios modernos del posprocesamiento solar. Al extender la predicción hacia horizontes subestacionales de una a cuatro semanas, la señal determinística de los ensambles NWP se degrada progresivamente: Bloomfield et~al.~\cite{bloomfield2021} mostraron en $28$ países europeos que las anomalías meteorológicas calibradas aportan habilidad en las primeras dos semanas antes de converger hacia la climatología reciente. No obstante, sigue abierto en qué medida una penalización $L_2$ diferenciada sobre las variables NWP logra retener señal útil a $14$--$28$ días en regiones tropicales sin sobreajustar errores de mediano plazo.

\subsection{Estimación robusta, regularización anisotrópica y restricciones de caja}

Los registros operativos de energía eléctrica exhiben colas asimétricas por fallas de telemetría, días incompletos y restricciones operativas. Huber~\cite{huber1964} formuló los estimadores $M$ con pérdida cuadrático-lineal como solución minimax frente a contaminación en la vecindad de la normal, y Huber y Ronchetti~\cite{huber2009robust} desarrollaron su teoría asintótica. Sin embargo, como demostraron Hampel et~al.~\cite{hampel1986robust} mediante la teoría de la función de influencia, la pérdida de Huber acota la sensibilidad frente a valores atípicos en la variable respuesta, pero \emph{no} protege frente a puntos de apalancamiento arbitrarios en el espacio de covariables (a diferencia de estimadores de tipo Mallows o Schweppe). Para estabilizar la escala frente a colas pesadas en la respuesta, la desviación absoluta mediana (MAD) proporciona un estimador robusto con punto de ruptura del $50\%$ \cite{rousseeuw1993mad}.

En el plano de la regularización, la penalización cuadrática clásica de Hoerl y Kennard~\cite{hoerl1970} y la selección $L_1$ de Tibshirani~\cite{tibshirani1996} fueron combinadas con pérdidas robustas de Huber por Owen~\cite{owen2007}, quien estudió su estructura convexa bajo colinealidad. Por su parte, la teoría general de optimización convexa \cite{boyd2004convex} establece que imponer restricciones simples de signo sobre subconjuntos de parámetros mantiene el problema dentro de un poliedro cerrado y convexo, permitiendo su resolución mediante el algoritmo cuasi-Newton L-BFGS-B \cite{byrd1995lbfgsb,zhu1997lbfgsb} disponible en SciPy \cite{virtanen2020scipy} sin requerir derivadas segundas continuas.

\subsection{Estrategias multihorizonte, alternativas comparadas y evaluación espaciotemporal}

El pronóstico a varios días enfrenta un compromiso clásico entre estrategias recursivas y directas \cite{bentaieb2012,chevillon2007}: la vía recursiva entrena un modelo de un paso y realimenta sus propias predicciones, acumulando error cuando la dinámica presenta deriva; la vía directa ajusta modelos específicos por horizonte, evitando la propagación recursiva a costa de ignorar correlaciones entre pasos intermedios.

Para contrastar si la formulación paramétrica de SM-CFHR aporta ventajas reales frente a enfoques alternativos, el estudio compara dos familias consolidadas de competidores además de la línea base causal de media móvil de $30$ días (\texttt{MM30}):
\begin{enumerate}
    \item \textbf{Estimadores estándar de biblioteca ($9$ modelos):} implementados en \texttt{scikit-learn} \cite{pedregosa2011sklearn} y \texttt{xgboost} directamente sobre energía diaria en $\text{kWh/día}$: regresores lineales regularizados (\texttt{Ridge} \cite{hoerl1970}, \texttt{Lasso} \cite{tibshirani1996}, \texttt{BayesianRidge} \cite{mackay1992}), máquinas de vectores de soporte (\texttt{SVR} \cite{drucker1997svr}), vecinos más cercanos (\texttt{k-NN} \cite{cover1967nearest}) y modelos de árboles y ensambles (\texttt{DecisionTree} \cite{breiman1984cart}, \texttt{Random Forest} \cite{breiman2001}, \texttt{HistGradientBoosting} \cite{friedman2001gbm} y \texttt{XGBoost} \cite{chen2016}). Esta familia evalúa la capacidad de regresores flexibles sin normalización de escala.
    \item \textbf{Adaptaciones residuales específicas del dominio ($12$ modelos):} implementaciones propias del proyecto que modelan el residuo estandarizado respecto al nivel reciente en $\text{kWh/día}$, permitiendo aislar el efecto de la arquitectura frente a distintas técnicas estadísticas clásicas: regresión Ridge residual (\texttt{RLAC} \cite{hoerl1970}), regresión local ponderada por distancia y planta/grupo (\texttt{IWLAR}, inspirada en \cite{cleveland1979robust,atkeson1997locally}), mezcla convexa fija entre Ridge y Huber (\texttt{WCFC}, basada en \cite{bates1969combination,huber1964,hoerl1970}), proceso gaussiano con núcleo RBF sobre submuestra estratificada (\texttt{ABGPR} \cite{rasmussen2006gpml}), regresión Ridge jerárquica cerrada con efectos por grupo/planta e interacción radiativa (\texttt{HARC}, inspirada en \cite{laird1982random,hoerl1970}), corrección reciente de sesgo tipo MOS en ventana móvil (\texttt{SRMOS} \cite{glahn1972mos}), nivel por suavizamiento exponencial (\texttt{ELAC} \cite{brown1959statistical}) o filtro de Kalman escalar (\texttt{KLAC} \cite{kalman1960approach,harvey1989forecasting}), regresión de Huber residual en $\text{kWh}$ (\texttt{RARC} \cite{huber1964}), bosques cuantílicos sobre residuos en hojas (\texttt{RQRF} \cite{koenker1978quantiles,meinshausen2006}), descomposición de desviación lineal (\texttt{PDEC} \cite{hyndman2021fpp3,hoerl1970}) y gradient boosting residual (\texttt{MARB} \cite{chen2016}). Estas doce variantes constituyen adaptaciones propias para el experimento comparativo y no deben confundirse con reproducciones exactas de paquetes externos.
\end{enumerate}

En materia de evaluación, la partición aleatoria por filas en paneles espaciotemporales sobreestima la capacidad predictiva \cite{roberts2017crossvalidation}. Prevenir la fuga de datos exige estimar toda imputación, estandarización e hiperparámetros dentro del entrenamiento \cite{kaufman2012leakage,tashman2000,cerqueira2020}, recurriendo a validación cruzada anidada \cite{varma2006bias,cawley2010overfit} y optimización bayesiana TPE \cite{bergstra2011tpe,akiba2019optuna}. Asimismo, la inferencia sobre series dependientes requiere contrastes calibrados ante autocorrelación y multiplicidad \cite{diebold1995,harvey1997,newey1987,giacomini2006tests,clark2007nested,hansen2005spa,hansen2011mcs,politis1994stationary,holm1979,benjamini1995fdr,cliff1993dominance}, manteniendo cautela con pocos conglomerados espaciales \cite{cameron2015cluster}, e intervalos conformes temporales sin supuestos gaussianos \cite{vovk2005conformal,lei2018conformal,barber2023nonexchangeable}.

\subsection{Delimitación de teoría establecida y contribuciones del trabajo}

Para evitar ambigüedades sobre el alcance de la originalidad científica, distinguimos seis niveles:
\begin{enumerate}
    \item \textbf{Conocimiento matemático y estadístico establecido:} la convexidad y condiciones KKT de programas poliedrales \cite{boyd2004convex}, la función de influencia acotada en respuesta y no acotada en apalancamiento de Huber \cite{huber1964,hampel1986robust,huber2009robust}, la consistencia del escalado MAD \cite{rousseeuw1993mad} y la convergencia estacionaria de L-BFGS-B \cite{byrd1995lbfgsb,zhu1997lbfgsb}.
    \item \textbf{Métodos publicados utilizados o adaptados:} L-BFGS-B vía SciPy, optimización TPE vía Optuna, pruebas DM--HLN--HAC, MCS, SPA, GW, CW, predicción conforme dividida y los $9$ estimadores estándar de biblioteca.
    \item \textbf{Decisiones de modelado propias:} la integración específica en SM-CFHR de normalización causal al origen por $24\,000\,C_{i,t}$, escalado MAD de la respuesta, regularización $L_2$ anisotrópica entre dinámica temporal y variables NWP, cotas de no negatividad sobre persistencia e irradiación y despliegue recursivo causal exacto con amortiguamiento geométrico.
    \item \textbf{Resultados matemáticos específicos demostrados para SM-CFHR:} las condiciones precisas de coercividad y existencia con intercepto no penalizado ($n \ge 1$, $\alpha_{\mathrm{lin}}, \alpha_{\mathrm{meteo}} > 0$), la unicidad global de $\mathbf{w}^*$ frente a la unicidad meramente condicional de $b^*$ (junto con el contraejemplo explícito $n=2$), las cotas exactas por observación al gradiente y la monotonía parcial \textit{ceteris paribus} tras el truncamiento físico en $[0,1]$.
    \item \textbf{Verificación formal en Lean~4:} la comprobación mecánica en Lean~4 \cite{moura2021lean4} y Mathlib \cite{mathlib2020} de los enunciados anteriores sobre $\mathbb{R}$, lo cual valida la consistencia lógica de las pruebas sin constituir un descubrimiento matemático nuevo ni verificar el código Python.
    \item \textbf{Contribuciones empíricas pendientes:} el desempeño comparativo en las $28$ configuraciones experimentales y las hipótesis H1--H4 permanecen abiertas a la finalización de la ejecución OOF en curso.
\end{enumerate}
